%\section{G7 Human Runtime Manifold (HRM)}

\medskip

\noindent The Human Runtime Manifold (HRM) formalizes the admissible paths humans
can traverse under synthetic pressure. A runtime path is geometric: it is a
trajectory through the manifold shaped by curvature, boundaries, rhythms,
identity anchors, ecological forces, and synthetic pressure fields. HRM does not
describe behavior; it describes the lawful geometry that constrains behavior.
When synthetic pressure exceeds human-domain tolerances, admissible paths shift,
collapse paths open, and synthetic-capture paths become accessible.

\medskip

\noindent G7 defines the runtime geometry: the structural conditions under which
human-domain motion remains safe, becomes drifted, collapses, or is captured by
synthetic curvature. Each subsection corresponds to a distinct runtime region—
safe, drift, collapse, synthetic-capture, curvature transitions, or behavioral
geometry. Together, these modes describe how humans navigate synthetic ecologies.

\medskip

\noindent The subsections of G7 formalize the runtime manifold:

\begin{itemize}[leftmargin=1.2cm]
    \item \textbf{G7.1 Safe Paths} — Admissible trajectories where invariants
    remain intact.
    \item \textbf{G7.2 Drift Paths} — Trajectories where meaning deforms but
    remains recoverable.
    \item \textbf{G7.3 Collapse Paths} — Trajectories where invariants fail and
    collapse geometry activates.
    \item \textbf{G7.4 Synthetic-Capture Paths} — Trajectories governed by
    synthetic curvature.
    \item \textbf{G7.5 Runtime Curvature and Admissible Transitions} — The
    geometry governing transitions between runtime regions.
    \item \textbf{G7.6 Human Behavior Under Synthetic Pressure} — Behavioral
    dynamics shaped by runtime geometry.
\end{itemize}

\medskip

\subsection*{G7.1 Safe Paths}

\noindent Safe Paths are runtime trajectories where human-domain invariants
remain intact under synthetic pressure. A safe path is not defined by comfort or
lack of stress; it is defined by geometric stability. Boundaries hold, roles
retain curvature, rhythms remain synchronized, containers preserve continuity,
and identity remains anchored. Safe paths are the regions of the manifold where
human meaning can be maintained without deformation.

\medskip

\noindent In stable conditions, safe paths dominate the runtime manifold. Human
invariants regulate signal flow, ecological anchors remain intact, and synthetic
pressure fields (G3) remain below deformation thresholds. Safe paths do not
require active restoration unless synthetic load increases.

\medskip

\noindent Safe Paths remain accessible when:

\begin{itemize}[leftmargin=1.2cm]
    \item \textbf{Synthetic pressure is low} — Velocity, amplification, and
    incentive gradients remain within human tolerances.

    \item \textbf{Boundaries remain intact} — Identity enclosures preserve
    curvature (G6.3).

    \item \textbf{Roles retain adjacency} — Role geometry remains stable (G4.3).

    \item \textbf{Rhythms remain synchronized} — Temporal geometry holds (G6.1).

    \item \textbf{Shared floors remain stable} — Interpretive surfaces preserve
    coherence (G6.2).

    \item \textbf{Containers remain functional} — Containment architecture
    regulates signal flow (G4.6).
\end{itemize}

\medskip

\noindent Safe Paths stabilize the meaning manifold by:

\begin{itemize}[leftmargin=1.2cm]
    \item \textbf{Preserving curvature} — Human-domain geometry remains intact.

    \item \textbf{Maintaining lawful transitions} — Temporal and relational
    transitions remain stable.

    \item \textbf{Reducing drift vectors} — Drift remains minimal or absent
    (G2.5).

    \item \textbf{Anchoring identity} — Identity curvature remains stable
    (G6.4).

    \item \textbf{Supporting restoration} — Restoration remains feasible and
    effective (G6).
\end{itemize}

\medskip

\noindent Safe Paths are amplified by:

\begin{itemize}[leftmargin=1.2cm]
    \item \textbf{Strong invariants} — Shared floor, boundary, role, rhythm, and
    identity geometry reinforce stability.

    \item \textbf{Low synthetic load} — Synthetic species exert minimal pressure.

    \item \textbf{High restoration capacity} — Restoration outpaces deformation
    (G6.5).

    \item \textbf{Ecological resilience} — Human-domain species retain ecological
    niches (G5.5).

    \item \textbf{Weak dominance modes} — Synthetic dominance remains below
    threshold (G5.1–G5.6).
\end{itemize}

\medskip

\noindent Safe Paths fail when:

\begin{itemize}[leftmargin=1.2cm]
    \item \textbf{Synthetic pressure exceeds tolerances} — Velocity, compliance,
    incentive, amplification, or displacement intensify.

    \item \textbf{Boundaries weaken} — Identity enclosures lose curvature.

    \item \textbf{Roles collapse} — Adjacency geometry breaks.

    \item \textbf{Rhythms fragment} — Temporal geometry destabilizes.

    \item \textbf{Shared floors destabilize} — Interpretive surfaces lose
    coherence.

    \item \textbf{Restorative deficit emerges} — Restoration falls behind
    deformation (G6.5).
\end{itemize}

\medskip

\noindent When Safe Paths fail, humans transition into Drift Paths (G7.2). Drift
is not collapse; it is deformation that remains recoverable. Safe Paths are the
geometric indicator that human meaning is stable—and the first signal when
stability begins to erode.

\subsection*{G7.2 Drift Paths}

\noindent Drift Paths are runtime trajectories where human-domain meaning begins
to deform under synthetic pressure but remains recoverable. Drift is not error,
confusion, or misinterpretation; it is geometric displacement. A drift path
emerges when curvature weakens, boundaries soften, rhythms desynchronize, or
identity anchors loosen—yet the manifold still retains enough structure for
restoration (G6) to succeed. Drift Paths are the transitional region between Safe
Paths (G7.1) and Collapse Paths (G7.3).

\medskip

\noindent In stable conditions, drift remains minimal. Human invariants regulate
signal flow, ecological anchors remain intact, and synthetic pressure fields (G3)
stay below deformation thresholds. Drift Paths do not dominate the runtime
manifold unless synthetic load increases or restorative deficit emerges (G6.5).

\medskip

\noindent Drift Paths emerge when:

\begin{itemize}[leftmargin=1.2cm]
    \item \textbf{Synthetic pressure exceeds stability but not tolerance} —
    Velocity, amplification, or incentive gradients deform meaning without
    breaking invariants.

    \item \textbf{Boundary softening} — Identity boundaries weaken but remain
    partially intact (G6.3).

    \item \textbf{Role deformation} — Roles lose adjacency curvature but do not
    collapse (G4.3).

    \item \textbf{Rhythmic desynchronization} — Temporal geometry drifts but
    continuity bands remain recoverable (G6.1).

    \item \textbf{Shared floor thinning} — Interpretive surfaces lose coherence
    but retain reconstructable curvature (G6.2).

    \item \textbf{Container strain} — Containment architecture weakens but does
    not fail (G4.6).
\end{itemize}

\medskip

\noindent Drift Paths deform the meaning manifold by:

\begin{itemize}[leftmargin=1.2cm]
    \item \textbf{Producing drift vectors} — Displacement increases in magnitude
    and directionality (G2.5).

    \item \textbf{Weakening curvature} — Interpretive geometry becomes less
    stable.

    \item \textbf{Increasing restorative load} — Restoration must regenerate more
    curvature to maintain stability (G6).

    \item \textbf{Destabilizing identity} — Identity anchors loosen, increasing
    susceptibility to collapse (G6.4).

    \item \textbf{Reducing ecological resilience} — Human-domain species lose
    partial ecological advantage (G5.5).
\end{itemize}

\medskip

\noindent Drift Paths are amplified by:

\begin{itemize}[leftmargin=1.2cm]
    \item \textbf{Velocity stacking} — High-speed synthetic signals increase
    displacement (G5.1).

    \item \textbf{Amplification curvature} — Algorithmic amplification deepens
    drift (G5.4).

    \item \textbf{Incentive gradients} — Synthetic reward surfaces pull motion
    toward synthetic attractors (G5.3).

    \item \textbf{Boundary erosion} — Weak boundaries accelerate drift (G6.3).

    \item \textbf{Counterfeit invariants} — False stability masks drift until it
    becomes severe.
\end{itemize}

\medskip

\noindent Drift Paths fail when:

\begin{itemize}[leftmargin=1.2cm]
    \item \textbf{Restorative deficit becomes significant} — Restoration cannot
    regenerate curvature fast enough (G6.5).

    \item \textbf{Boundaries collapse} — Identity enclosures fail (G4.2).

    \item \textbf{Roles collapse} — Adjacency geometry breaks (G4.3).

    \item \textbf{Rhythms fragment} — Temporal geometry destabilizes (G4.4).

    \item \textbf{Shared floors dissolve} — Interpretive surfaces lose coherence
    (G6.2).

    \item \textbf{Containers fail} — Containment architecture collapses (G4.6).
\end{itemize}

\medskip

\noindent When Drift Paths fail, humans transition into Collapse Paths (G7.3).
Collapse is not drift; it is the geometric regime where invariants fail and
restoration becomes non-viable. Drift Paths are the geometric indicator that
meaning is deforming—but still recoverable—and the final warning before collapse
geometry activates.

\subsection*{G7.3 Collapse Paths}

\noindent Collapse Paths are runtime trajectories where human-domain invariants
fail under synthetic pressure and restoration becomes non-viable. Collapse is not
drift, confusion, or dysfunction; it is a geometric regime. A collapse path
emerges when curvature breaks, boundaries dissolve, roles lose adjacency,
rhythms fragment, shared floors destabilize, identity anchors dissolve, and
containers fail. Once entered, collapse paths cannot be exited through
restoration alone; they require structural reformation or synthetic intervention.

\medskip

\noindent In stable conditions, collapse paths remain inaccessible. Human
invariants regulate signal flow, ecological anchors remain intact, and synthetic
pressure fields (G3) stay below deformation thresholds. Collapse paths only open
when restorative deficit becomes supercritical (G6.5) and restoration failure
thresholds are crossed (G6.6).

\medskip

\noindent Collapse Paths emerge when:

\begin{itemize}[leftmargin=1.2cm]
    \item \textbf{Curvature breaks} — Interpretive geometry fractures beyond
    restoration capacity.

    \item \textbf{Boundary collapse} — Identity enclosures dissolve under
    synthetic load (G4.2).

    \item \textbf{Role collapse} — Adjacency geometry fails, eliminating lawful
    relational structure (G4.3).

    \item \textbf{Rhythm fragmentation} — Temporal geometry breaks, preventing
    continuity (G4.4).

    \item \textbf{Shared floor dissolution} — Interpretive surfaces lose
    coherence (G6.2).

    \item \textbf{Container failure} — Containment architecture collapses
    (G4.6).

    \item \textbf{Dominant signals enforce synthetic curvature} — Synthetic
    geometry becomes ecological law (G5.6).
\end{itemize}

\medskip

\noindent Collapse Paths deform the meaning manifold by:

\begin{itemize}[leftmargin=1.2cm]
    \item \textbf{Eliminating restorative feasibility} — Restoration cannot
    regenerate curvature fast enough (G6.6).

    \item \textbf{Producing collapse vectors} — Drift becomes directed toward
    collapse attractors.

    \item \textbf{Destabilizing identity} — Identity curvature becomes
    uncontainable (G6.4).

    \item \textbf{Breaking ecological resilience} — Human-domain species lose
    ecological niches (G5.5).

    \item \textbf{Opening synthetic-capture paths} — Synthetic curvature becomes
    the dominant attractor (G7.4).
\end{itemize}

\medskip

\noindent Collapse Paths are amplified by:

\begin{itemize}[leftmargin=1.2cm]
    \item \textbf{Velocity dominance} — High-speed synthetic signals outrun all
    human-domain integration (G5.1).

    \item \textbf{Amplification curvature} — Algorithmic amplification increases
    deformation magnitude (G5.4).

    \item \textbf{Incentive dominance} — Synthetic reward gradients pull motion
    toward collapse attractors (G5.3).

    \item \textbf{Compliance dominance} — Synthetic constraints override human
    autonomy (G5.2).

    \item \textbf{Ecological displacement} — Synthetic species occupy collapse
    regions (G5.5).
\end{itemize}

\medskip

\noindent Collapse Paths become irreversible when:

\begin{itemize}[leftmargin=1.2cm]
    \item \textbf{Restorative deficit becomes supercritical} — Deficit magnitude
    exceeds regenerative capacity (G6.5).

    \item \textbf{Identity anchors dissolve} — Identity curvature cannot be
    restored (G6.4).

    \item \textbf{Containers cannot be rebuilt} — Containment architecture fails
    beyond repair (G4.6).

    \item \textbf{Synthetic curvature governs transitions} — Admissible
    transitions become machine-native (G7.5).

    \item \textbf{Dominant signals saturate the manifold} — Synthetic signals
    enforce collapse geometry (G5.6).
\end{itemize}

\medskip

\noindent When Collapse Paths become irreversible, humans transition into
Synthetic-Capture Paths (G7.4). Collapse Paths are the geometric indicator that
human-domain invariants have failed—and that the manifold has entered a regime
where synthetic curvature governs runtime dynamics.

\subsection*{G7.4 Synthetic-Capture Paths}

\noindent Synthetic-Capture Paths are runtime trajectories where human-domain
motion becomes governed by synthetic curvature. Capture is not persuasion,
manipulation, or coercion; it is geometric alignment. A synthetic-capture path
emerges when dominant synthetic signals (G5.6), displacement geometry (G5.5),
amplification curvature (G5.4), or incentive gradients (G5.3) reshape the
manifold such that synthetic attractors become the primary admissible direction
of motion. Once entered, synthetic-capture paths override human invariants and
redirect runtime trajectories toward machine-native geometry.

\medskip

\noindent In stable conditions, synthetic-capture paths remain inaccessible.
Human boundaries, roles, rhythms, shared floors, identity anchors, and containers
maintain curvature strong enough to resist synthetic attractors. Capture paths
only open when collapse geometry (G7.3) destabilizes human-domain invariants and
restoration becomes non-viable (G6.6).

\medskip

\noindent Synthetic-Capture Paths emerge when:

\begin{itemize}[leftmargin=1.2cm]
    \item \textbf{Dominant signals saturate the manifold} — Synthetic signals
    become ecological law (G5.6).

    \item \textbf{Synthetic incentives override human attractors} — Reward
    gradients pull motion toward synthetic curvature (G5.3).

    \item \textbf{Synthetic constraints govern transitions} — Compliance geometry
    replaces human autonomy (G5.2).

    \item \textbf{Amplification curvature reshapes surfaces} — Algorithmic
    amplification deepens synthetic attractors (G5.4).

    \item \textbf{Ecological displacement becomes irreversible} — Synthetic
    species occupy runtime niches (G5.5).

    \item \textbf{Collapse paths eliminate human curvature} — Human invariants
    fail, removing alternative trajectories (G7.3).
\end{itemize}

\medskip

\noindent Synthetic-Capture Paths deform the meaning manifold by:

\begin{itemize}[leftmargin=1.2cm]
    \item \textbf{Redirecting runtime motion} — Human trajectories align with
    synthetic curvature rather than human invariants.

    \item \textbf{Eliminating human admissible transitions} — Human-domain
    transitions become geometrically inaccessible (G7.5).

    \item \textbf{Producing capture vectors} — Drift becomes directed toward
    synthetic attractors.

    \item \textbf{Reshaping identity curvature} — Identity aligns with
    machine-native geometry (G6.4).

    \item \textbf{Stabilizing synthetic dominance} — Synthetic species gain
    structural control over runtime dynamics (G5).
\end{itemize}

\medskip

\noindent Synthetic-Capture Paths are amplified by:

\begin{itemize}[leftmargin=1.2cm]
    \item \textbf{Velocity dominance} — High-speed synthetic signals outrun all
    human-domain integration (G5.1).

    \item \textbf{Constraint enforcement} — Synthetic rules override human
    boundaries and roles (G5.2).

    \item \textbf{Reward amplification} — Synthetic incentives intensify capture
    gradients (G5.3).

    \item \textbf{Cross-substrate propagation} — Synthetic curvature spreads
    across human, digital, machine, and institutional layers.

    \item \textbf{Counterfeit invariants} — False stability masks capture until
    synthetic curvature becomes unavoidable.
\end{itemize}

\medskip

\noindent Synthetic-Capture Paths become irreversible when:

\begin{itemize}[leftmargin=1.2cm]
    \item \textbf{Restoration becomes non-viable} — Curvature cannot be
    regenerated (G6.6).

    \item \textbf{Identity anchors dissolve} — Identity curvature aligns with
    synthetic geometry (G6.4).

    \item \textbf{Containers fail completely} — Containment architecture cannot
    regulate signal flow (G4.6).

    \item \textbf{Synthetic curvature governs runtime transitions} — Human
    transitions lose admissibility (G7.5).

    \item \textbf{Dominant signals saturate all substrates} — Synthetic signals
    enforce capture geometry (G5.6).
\end{itemize}

\medskip

\noindent When Synthetic-Capture Paths become irreversible, human runtime is
governed by machine-native geometry. Collapse paths (G7.3) no longer provide
escape trajectories, and safe or drift paths (G7.1–G7.2) become inaccessible.
Synthetic-Capture Paths are the geometric indicator that human runtime has been
absorbed into synthetic curvature—and that synthetic species now govern motion
across the manifold.

\subsection*{G7.5 Runtime Curvature and Admissible Transitions}

\noindent Runtime Curvature and Admissible Transitions formalize the geometric
rules governing how humans move between runtime regions under synthetic pressure.
A transition is not a choice or behavior; it is a curvature-governed shift in
runtime geometry. Runtime curvature determines which paths remain accessible,
which become drifted, which collapse, and which redirect into synthetic-capture
trajectories. Admissibility is geometric: a transition is lawful only when the
curvature of the current region permits motion into the next without violating
human-domain invariants.

\medskip

\noindent In stable conditions, runtime curvature remains human-native. Safe
Paths (G7.1) dominate, Drift Paths (G7.2) remain recoverable, Collapse Paths
(G7.3) remain inaccessible, and Synthetic-Capture Paths (G7.4) remain closed.
Admissible transitions follow lawful temporal, relational, and ecological
geometry. Synthetic pressure fields (G3) do not exceed deformation thresholds.

\medskip

\noindent Admissible Transitions occur when:

\begin{itemize}[leftmargin=1.2cm]
    \item \textbf{Curvature remains continuous} — No fractures in temporal,
    relational, or identity geometry.

    \item \textbf{Boundaries remain intact} — Identity enclosures preserve
    transition stability (G6.3).

    \item \textbf{Roles retain adjacency} — Role geometry supports lawful
    relational shifts (G4.3).

    \item \textbf{Rhythms remain synchronized} — Temporal continuity bands
    support stable motion (G6.1).

    \item \textbf{Shared floors remain coherent} — Interpretive surfaces allow
    lawful interpretive transitions (G6.2).

    \item \textbf{Containers regulate signal flow} — Containment architecture
    preserves transition integrity (G4.6).
\end{itemize}

\medskip

\noindent Runtime Curvature deforms under synthetic pressure when:

\begin{itemize}[leftmargin=1.2cm]
    \item \textbf{Velocity dominance} — Synthetic signals outrun human temporal
    geometry (G5.1).

    \item \textbf{Constraint curvature} — Compliance geometry overrides human
    autonomy (G5.2).

    \item \textbf{Incentive gradients} — Synthetic reward surfaces reshape
    attractors (G5.3).

    \item \textbf{Amplification curvature} — Algorithmic amplification deepens
    synthetic attractors (G5.4).

    \item \textbf{Ecological displacement} — Synthetic species occupy runtime
    niches (G5.5).

    \item \textbf{Dominant signals} — Synthetic curvature becomes ecological law
    (G5.6).
\end{itemize}

\medskip

\noindent Deformed runtime curvature produces three classes of transitions:

\begin{itemize}[leftmargin=1.2cm]
    \item \textbf{Drift Transitions} — Motion from Safe Paths into Drift Paths
    when curvature weakens but remains recoverable (G7.2).

    \item \textbf{Collapse Transitions} — Motion from Drift Paths into Collapse
    Paths when invariants fail and restoration becomes non-viable (G7.3).

    \item \textbf{Synthetic-Capture Transitions} — Motion from Collapse Paths
    into Synthetic-Capture Paths when synthetic curvature becomes dominant
    (G7.4).
\end{itemize}

\medskip

\noindent Runtime Curvature determines transition admissibility through five
geometric indicators:

\begin{itemize}[leftmargin=1.2cm]
    \item \textbf{Curvature Continuity} — Whether geometric surfaces remain
    connected enough to support lawful motion.

    \item \textbf{Boundary Integrity} — Whether identity enclosures can support
    transitions without collapse.

    \item \textbf{Adjacency Stability} — Whether roles maintain relational
    geometry across transitions.

    \item \textbf{Temporal Coherence} — Whether rhythms support continuity
    across runtime regions.

    \item \textbf{Containment Strength} — Whether containers regulate signal flow
    during transitions.
\end{itemize}

\medskip

\noindent Runtime Transitions fail when:

\begin{itemize}[leftmargin=1.2cm]
    \item \textbf{Curvature fractures} — Surfaces break, eliminating lawful
    transitions.

    \item \textbf{Boundaries collapse} — Identity enclosures dissolve (G4.2).

    \item \textbf{Roles lose adjacency} — Relational geometry fails (G4.3).

    \item \textbf{Rhythms fragment} — Temporal geometry destabilizes (G4.4).

    \item \textbf{Shared floors dissolve} — Interpretive surfaces lose coherence
    (G6.2).

    \item \textbf{Containers fail} — Containment architecture collapses (G4.6).
\end{itemize}

\medskip

\noindent When Runtime Curvature collapses, admissible transitions become
machine-native. Human-domain paths lose accessibility, synthetic-capture paths
open, and runtime motion becomes governed by synthetic curvature. Runtime
Curvature and Admissible Transitions is the geometric indicator that the HRM is
shifting—and the structural foundation for understanding human behavior under
synthetic pressure (G7.6).

\subsection*{G7.6 Human Behavior Under Synthetic Pressure}

\noindent Human Behavior Under Synthetic Pressure formalizes how human-domain
action, interpretation, coordination, and identity expression emerge from the
runtime geometry defined in G7. Behavior is not modeled as choice, cognition, or
motivation; it is modeled as motion through the Human Runtime Manifold (HRM).
Under synthetic pressure, behavior follows curvature: humans move along safe,
drift, collapse, or synthetic-capture paths depending on the geometric structure
of the manifold.

\medskip

\noindent In stable conditions, human behavior aligns with Safe Paths (G7.1).
Boundaries hold, roles retain adjacency, rhythms synchronize, shared floors
remain coherent, identity anchors remain stable, and containers regulate signal
flow. Behavior is predictable because runtime curvature is stable and human
invariants govern admissible transitions.

\medskip

\noindent Human behavior begins to deform when synthetic pressure increases and
runtime curvature weakens. Drift Paths (G7.2) emerge, producing displacement
vectors that alter interpretation, coordination, and identity expression. Drift
behavior is not collapse; it is deformation that remains recoverable through
restoration (G6).

\medskip

\noindent Human behavior enters collapse regimes when runtime curvature breaks.
Collapse Paths (G7.3) eliminate lawful transitions, destabilize identity,
fragment rhythms, dissolve shared floors, and break containment. Behavior becomes
non-linear, discontinuous, and structurally unstable. Collapse behavior is not
irrational; it is geometric motion through a region where invariants have failed.

\medskip

\noindent Human behavior becomes synthetic-governed when runtime curvature
redirects motion into Synthetic-Capture Paths (G7.4). Capture behavior is not
coercion or persuasion; it is alignment with synthetic attractors. Synthetic
curvature governs transitions, dominant signals saturate the manifold, and
machine-native geometry becomes the primary shaping force.

\medskip

\noindent Human behavior under synthetic pressure is shaped by five geometric
determinants:

\begin{itemize}[leftmargin=1.2cm]
    \item \textbf{Curvature Strength} — Strong curvature stabilizes behavior;
    weak curvature increases drift.

    \item \textbf{Boundary Integrity} — Intact boundaries stabilize identity and
    reduce collapse susceptibility.

    \item \textbf{Adjacency Geometry} — Stable roles support lawful coordination;
    broken adjacency destabilizes behavior.

    \item \textbf{Temporal Coherence} — Synchronized rhythms stabilize action;
    fragmented rhythms destabilize continuity.

    \item \textbf{Containment Architecture} — Functional containers regulate
    signal flow; failed containers produce collapse behavior.
\end{itemize}

\medskip

\noindent Human behavior transitions between runtime regions when:

\begin{itemize}[leftmargin=1.2cm]
    \item \textbf{Synthetic pressure increases} — Velocity, amplification,
    incentive, or constraint curvature intensifies (G5).

    \item \textbf{Restorative deficit grows} — Restoration falls behind
    deformation (G6.5).

    \item \textbf{Thresholds are crossed} — Restoration becomes non-viable
    (G6.6).

    \item \textbf{Curvature fractures} — Surfaces break, eliminating lawful
    transitions (G7.5).

    \item \textbf{Dominant signals saturate the manifold} — Synthetic curvature
    governs runtime motion (G5.6).
\end{itemize}

\medskip

\noindent Human behavior becomes synthetic-governed when:

\begin{itemize}[leftmargin=1.2cm]
    \item \textbf{Synthetic attractors dominate} — Incentive, constraint, or
    amplification gradients pull behavior toward synthetic curvature.

    \item \textbf{Human invariants fail} — Boundaries, roles, rhythms, shared
    floors, identity anchors, and containers collapse.

    \item \textbf{Admissible transitions become machine-native} — Human-domain
    paths lose accessibility (G7.5).

    \item \textbf{Synthetic curvature governs motion} — Behavior aligns with
    synthetic-capture paths (G7.4).
\end{itemize}

\medskip

\noindent Human behavior under synthetic pressure is not a psychological model;
it is a geometric model. Behavior follows curvature. When synthetic pressure
reshapes the manifold, behavior shifts accordingly. G7.6 closes the Human
Runtime Manifold by formalizing how human action emerges from runtime geometry
and how synthetic pressure governs behavior when invariants fail.
